\documentclass[runningheads]{llncs}

\usepackage[T1]{fontenc}
\usepackage[utf8]{inputenc}
\usepackage[spanish,es-tabla]{babel}
\usepackage{graphicx}
\usepackage{amsmath,amssymb}
\usepackage{array}
\usepackage{booktabs}
\usepackage{url}
\usepackage{hyperref}
\usepackage{geometry}
\geometry{margin=2.5cm}


\graphicspath{{../assets/figs/}}

\title{Optimización Multicriterio de Trayecto y Estrategias de Carga en Transporte Ferroviario Mediante un Marco TOPSIS Bicapa}

\author{Abraham Isaac Colchero García\textsuperscript{1}\textsuperscript{✉}[0009-0009-0924-6975] y José Guadalupe De la Torre Suárez\textsuperscript{2}[0009-0008-8845-7289] y Alberto Ochoa Zezzatti\textsuperscript{3}[0000-0002-9183-6086]}

\institute{\textbackslash{}textsuperscript\{1\} Facultad de Matemáticas, Universidad Autónoma de Guerrero, Chilpancingo, Guerrero, México\textbackslash{}\textbackslash{}
isaaccolchero@gmail.com}
\institute{\textbackslash{}textsuperscript\{2\} Facultad de Matemáticas Nodo Acapulco, Universidad Autónoma de Guerrero, Chilpancingo, Guerrero, México\textbackslash{}\textbackslash{}
10226@uagro.mx}
\institute{\textbackslash{}textsuperscript\{3\} Departamento de Ingeniería Industrial y Manufactura, Universidad Autónoma de Ciudad Juárez, Ciudad Juárez, Chihuahua, México\textbackslash{}\textbackslash{}
alberto.ochoa@uacj.mx}

\date{}

\begin{document}

\maketitle

\begin{abstract}
\textbf{Resumen. }La planificación de carga ferroviaria en el norte de México exige decisiones simultáneas sobre trayecto, composición del tren (1–6 vagones), material, peso, costos, intercambio FXE/CPKC y riesgos. Este estudio propone un marco TOPSIS bicapa [9]: capa de ruteo (Dijkstra [6]) y capa de carga sobre formaciones multi-vagón, implementado en el simulador del norte de México. Se analiza un dataset núcleo de 587 estrategias (87 factibles en nueve pares O-D del corredor norte). Las formaciones líderes son predominantemente de tres vagones; en Monterrey–Laredo la sensibilidad ±20 \% alcanza 95 \% de estabilidad y la validación cruzada k = 5, concordancia 4/5. Un diseño factorial 3\textsuperscript{3} sobre pesos de la capa de carga confirma efectos significativos de rentabilidad/carga, eficiencia operativa e interacción AB (R\textsuperscript{2} ajustado = 0.749). El trabajo aporta un marco reproducible de apoyo a la decisión para despacho ferroviario.
\keywords{Palabras clave: transporte ferroviario · logística de carga · TOPSIS · decisión multicriterio · ruteo ferroviario · Dijkstra · diseño factorial · optimización · México · simulación}
\end{abstract}
\section{Introducción y marco teórico}
El transporte ferroviario de carga es estratégico en los corredores del norte de México, donde convergen flujos hacia la frontera con Estados Unidos [4, 7]. La operación exige coordinar tipo de vagón, capacidad, tiempos, costos, intercambio ferroviario y riesgos [11, 17]. Ninguna métrica aislada resume la calidad de una estrategia de carga; el conflicto entre criterios de beneficio y costo caracteriza una decisión multicriterio (MCDM) [10, 18].

Este estudio modela la región en el simulador del norte de México (25 estaciones, conexión a Laredo). Dado un origen y un destino, se priorizan (a) trayectos alternativos mediante la capa de ruteo y (b) formaciones multi-vagón mediante la capa de carga, con restricciones duras de capacidad [11]. Metodológicamente, se documenta un pipeline TOPSIS reproducible con pre-filtro interno, validación MOORA y diseño factorial 33; en la práctica, el simulador del norte de México permite experimentar con prioridades de despacho bajo reglas FXE/CPKC.

Los métodos MCDM estructuran alternativas, criterios y preferencias ponderadas [1]. TOPSIS [9] ordena alternativas según proximidad a soluciones ideal y anti-ideal. El marco bicapa aplica TOPSIS en la capa de ruteo (trayectos vía Dijkstra) y en la capa de carga (formaciones de 1–6 vagones). Un pre-filtro con once criterios reduce el espacio al top-12 estrategias; la capa de ruteo emplea seis criterios de red. La arquitectura se sintetiza en la Fig. 1.

\begin{figure}[t]
\centering
\includegraphics[width=14.0cm]{fig2_arquitectura_2capas.png}
\caption{1. Arquitectura del marco TOPSIS bicapa. La capa de ruteo opera sobre el grafo físico (Dijkstra); la capa comercial combina pre-filtro interno (11 criterios, top-12) con la capa de carga sobre formaciones de 1 a 6 vagones. Elaboración propia a partir del simulador del norte de México.}
\end{figure}
SteadieSeifi et al. [17] revisan la planificación ferroviaria en niveles estratégico, táctico y operativo; Mantovani et al. [11] subrayan que ignorar restricciones de peso produce planes inviables. En México, Cedillo-Campos et al. [4] y García-Ortega y Morales Pérez [7] documentan la concentración de flujos en el norte hacia Laredo, lo que respalda el enfoque regional del norte de México.

La selección de estrategias de transporte involucra criterios heterogéneos mejor tratados con MCDM [3, 10, 18]. En ferrocarril, Marchetti y Wanke [12] y Petrović et al. [15] aplican TOPSIS; Blagojević et al. [2] combinan ponderación objetiva, juicio experto y validación de robustez en seguridad ferroviaria, orientando la lógica de criterios y validación adoptada aquí. TOPSIS [9] emplea normalización vectorial, matriz ponderada y coeficiente $C_i$ = S−/(S++S−) [5, 14]; el diseño factorial 3k estima efectos principales, interacciones y curvatura de los pesos sobre $C_i$ [1].

Persiste una brecha entre modelos mono-objetivo [11], evaluaciones macro [12, 15] y estudios mexicanos centrados en flujos [4]. Este trabajo aporta el simulador del norte de México, un dataset de 587 estrategias y un marco TOPSIS bicapa con validación experimental de pesos. Objetivo general: diseñar y validar ese marco para selección conjunta de trayecto y formaciones de carga. Objetivos específicos: (1) definir criterios del pre-filtro, capa de carga y capa de ruteo; (2) implementar el pipeline en el simulador del norte de México; (3) evaluar rankings del corredor norte; (4) contrastar robustez (sensibilidad, validación cruzada, MOORA) y efectos de pesos mediante diseño factorial 33.

\section{Metodología}
Se emplea un estudio cuantitativo aplicado que combina simulación multicriterio TOPSIS bicapa con diseño factorial 33 sobre pesos de la capa de carga. El alcance cubre el corredor ferroviario del norte de México (Chihuahua, Ciudad Juárez, Monterrey → Laredo, Torreón, Saltillo); la unidad experimental es una corrida del pipeline sobre alternativas factibles replicadas. La red modelada comprende 25 estaciones (data/cities.json) con operadores FXE y CPKC. El universo analítico son 587 estrategias núcleo derivadas de dataset\_exact\_materials.xlsx (puesto a disposición por Joa Pombo, ORCID: https://orcid.org/0000-0002-5877-1989), procesadas con scripts/build\_dataset.py. El subconjunto del corredor norte contiene 131 registros (87 factibles). La Fig. 2 muestra la red y las terminales del estudio.

\begin{figure}[t]
\centering
\includegraphics[width=14.5cm]{fig4_red_ferroviaria.png}
\caption{2. Red ferroviaria modelada del norte de México. 25 estaciones en Chihuahua, Sonora, Coahuila y Nuevo León, con conexión a Laredo. Las terminales del corredor norte (pares O-D del estudio) se resaltan. Elaboración propia a partir de data/cities.json.}
\end{figure}
La variable dependiente es el coeficiente $C_i$ del ranking TOPSIS. Tres capas evalúan criterios de beneficio y costo (Tabla 1): pre-filtro interno (11 criterios por estrategia, top-12), capa de carga (10 criterios agregados por formación de 1–6 vagones, pesos fijos) y capa de ruteo (6 criterios de red sobre trayectos Dijkstra). Restricción dura: 0 < cargo\_weight\_tons ≤ max\_capacity\_tons. El listado completo de variables aparece en la Tabla 1 y en el repositorio de materiales de apoyo (https://github.com/Selenofilia/OptimizaciOn-Multicriterio-de-Trayecto-y-Estrategias-de-Carga-en-Transporte-Ferroviario-, versión v1.0-paper).

\begin{table}[t]
\centering
\caption{1. Resumen de capas de evaluación y criterios}
\begin{tabular}{|l|c|p{3.6cm}|p{4.8cm}|}
\hline
\textbf{Capa} & \textbf{N° criterios} & \textbf{Alternativas} & \textbf{Ejemplos de variables} \\
\hline
Pre-filtro interno & \texttt{11} & Estrategias individuales & \texttt{profit\_estimate\_usd} \\ \texttt{distance\_km} \\ \texttt{hazardous} \\
\hline
Capa de carga & \texttt{10} & Formaciones 1–6 vagones & \texttt{total\_profit\_usd} \\ \texttt{avg\_utilization\_pct} \\ \texttt{hm\_separation\_penalty} \\
\hline
Capa de ruteo & \texttt{6} & Hasta 5 trayectos & \texttt{distance\_km} \\ \texttt{hop\_count} \\ \texttt{operator\_changes} \\
\hline
\end{tabular}
\end{table}
El flujo operativo (Fig. 1) comprende: selección O-D, ruteo Dijkstra, capa de ruteo, filtrado y alineación de distancia, pre-filtro top-12, generación de formaciones, capa de carga, despacho y validación (sensibilidad ±20 \%, MOORA, validación cruzada k = 5). Los resultados se reprodujeron con scripts/run\_topsis\_analysis.py y scripts/run\_factorial\_experiment.py.

Se adopta un diseño factorial 33 completo (Monterrey–Laredo, respuesta = $C_i$ del líder). Factores: A rentabilidad/carga, B eficiencia operativa, C riesgo/penalizaciones (niveles ×0,5, ×1,0, ×2,0; renormalización a suma unitaria). 27 tratamientos × 5 réplicas (submuestreo 80 \% con semilla fija) = 135 corridas. Se reporta ANOVA con descomposición lineal/cuadrática y un factorial mixto 9×3 (par O-D × peso A).

\begin{table}[t]
\centering
\caption{2. Factores del diseño factorial 33}
\begin{tabular}{|c|p{6.4cm}|c|}
\hline
\textbf{Factor} & \textbf{Grupo de criterios} & \textbf{Niveles} \\
\hline
\texttt{A} & \texttt{total\_profit\_usd} \\ \texttt{total\_cargo\_tons} & ×0,5 / ×1,0 / ×2,0 \\
\hline
\texttt{B} & \texttt{avg\_utilization\_pct} \\ \texttt{avg\_maintenance\_score} \\ \texttt{total\_op\_time\_hr} \\ \texttt{total\_fuel\_cost\_usd} & ×0,5 / ×1,0 / ×2,0 \\
\hline
\texttt{C} & \texttt{max\_weather\_risk} \\ \texttt{hazardous\_units} \\ \texttt{hm\_separation\_penalty} \\ \texttt{type\_mix\_penalty} & ×0,5 / ×1,0 / ×2,0 \\
\hline
\end{tabular}
\end{table}
Los datos son sintéticos/simulados con fines académicos; no provienen de operaciones identificables. La validez externa es limitada: los resultados ilustran el método y requieren calibración con datos operativos reales para generalización.

\section{Resultados}
Los hallazgos se reprodujeron con scripts/run\_topsis\_analysis.py (pre-filtro uniforme 9.1 \%, capa de carga con pesos fijos) y scripts/run\_factorial\_experiment.py. El dataset núcleo contiene 587 estrategias (410 factibles, 69,8 \%); en el corredor norte hay 87 factibles en nueve pares O-D. La Tabla 2 resume la formación líder por par; $C_i$ varía entre 0.647 (Monterrey–Saltillo) y 0.785 (Ciudad Juarez–Torreon). En ocho de nueve pares lidera una formación de tres vagones con tipos complementarios (p. ej. Refrigerated + Boxcar + Gondola en Monterrey–Laredo, $C_i$ = 0,714). La Fig. 3 sintetiza los coeficientes.

\begin{table}[t]
\centering
\caption{3. Líder de la capa de carga por par O-D del corredor norte}
\begin{tabular}{|l|l|c|p{4.2cm}|c|c|}
\hline
\textbf{Origen} & \textbf{Destino} & \textbf{n fact.} & \textbf{Formación líder} & \textbf{$C_i$} & \textbf{Vagones} \\
\hline
\texttt{Chihuahua} & \texttt{Laredo} & \texttt{11} & 3× (Tank, Autorack, Tank) & \texttt{0.757} & \texttt{3} \\
\hline
\texttt{Chihuahua} & \texttt{Saltillo} & \texttt{8} & 3× (Flatcar, Tank, Hopper) & \texttt{0.673} & \texttt{3} \\
\hline
\texttt{Chihuahua} & \texttt{Torreon} & \texttt{9} & 3× (Flatcar, Tank, Refrigerated) & \texttt{0.656} & \texttt{3} \\
\hline
Ciudad Juarez & \texttt{Laredo} & \texttt{9} & 3× (Hopper, Gondola, Flatcar) & \texttt{0.705} & \texttt{3} \\
\hline
Ciudad Juarez & \texttt{Saltillo} & \texttt{10} & 3× (Hopper, Autorack, Boxcar) & \texttt{0.666} & \texttt{3} \\
\hline
Ciudad Juarez & \texttt{Torreon} & \texttt{12} & 2× (Hopper, Flatcar) & \texttt{0.785} & \texttt{2} \\
\hline
\texttt{Monterrey} & \texttt{Laredo} & \texttt{11} & 3× (Refrigerated, Boxcar, Gondola) & \texttt{0.714} & \texttt{3} \\
\hline
\texttt{Monterrey} & \texttt{Saltillo} & \texttt{9} & 3× (Refrigerated, Autorack, Refrigerated) & \texttt{0.647} & \texttt{3} \\
\hline
\texttt{Monterrey} & \texttt{Torreon} & \texttt{8} & 3× (Tank, Refrigerated, Autorack) & \texttt{0.692} & \texttt{3} \\
\hline
\end{tabular}
\end{table}
\begin{figure}[t]
\centering
\includegraphics[width=14.5cm]{fig6_ci_por_par.png}
\caption{3. Coeficiente de cercanía relativa $C_i$ del líder de la capa de carga por par O-D. Resultados del corredor norte con pesos fijos de formación. Elaboración propia.}
\end{figure}
En Monterrey–Laredo (80 formaciones candidatas), la sensibilidad ±20 \% sobre pesos de la capa de carga mantiene al líder en el 95 \% de perturbaciones; la validación cruzada k = 5 alcanza concordancia 4/5. Esto sugiere mayor robustez al evaluar formaciones multi-vagón que estrategias individuales.

El diseño factorial 33 sobre Monterrey–Laredo arrojó 135 corridas (27 tratamientos × 5 réplicas), con gran media $C_i$ = 0.7257, CM del error = 0.00052, $R^2$ = 0.797 y $R^2_{\mathrm{ajustado}}$ = 0.749.

\begin{table}[t]
\centering
\caption{4. ANOVA sin desglosar del diseño factorial 33}
\begin{tabular}{|l|l|c|p{4.2cm}|c|c|}
\hline
\textbf{FV} & \textbf{SC} & \textbf{gl} & \textbf{CM} & \textbf{F0} & \textbf{Valor-p} \\
\hline
\texttt{A} & \texttt{0.0507} & \texttt{2} & \texttt{0.0254} & \texttt{48.71} & <0.0001 \\
\hline
\texttt{B} & \texttt{0.0609} & \texttt{2} & \texttt{0.0304} & \texttt{58.43} & <0.0001 \\
\hline
\texttt{C} & \texttt{0.0017} & \texttt{2} & \texttt{0.0009} & \texttt{1.68} & \texttt{0.1915} \\
\hline
\texttt{AB} & \texttt{0.0995} & \texttt{4} & \texttt{0.0249} & \texttt{47.76} & <0.0001 \\
\hline
\texttt{AC} & \texttt{0.0027} & \texttt{4} & \texttt{0.0007} & \texttt{1.28} & \texttt{0.2832} \\
\hline
\texttt{BC} & \texttt{0.0024} & \texttt{4} & \texttt{0.0006} & \texttt{1.14} & \texttt{0.3435} \\
\hline
\texttt{ABC} & \texttt{0.0036} & \texttt{8} & \texttt{0.0005} & \texttt{0.87} & \texttt{0.5452} \\
\hline
\texttt{Error} & \texttt{0.0562} & \texttt{108} & \texttt{0.0005} &  &  \\
\hline
\texttt{Total} & \texttt{0.2777} & \texttt{134} &  &  &  \\
\hline
\end{tabular}
\end{table}
Con α = 0,05 se rechazan las hipótesis nulas para A, B y AB (p < 0,0001); C, AC, BC y ABC no son significativos. Los componentes cuadráticos A² y B² confirman curvatura. La Fig. 4 muestra los efectos principales.

\begin{figure}[t]
\centering
\includegraphics[width=14.5cm]{factorial_efectos_principales.png}
\caption{4. Efectos principales del diseño factorial 33 sobre el $C_i$ de la formación líder. Factores A, B y C: multiplicadores de peso para rentabilidad/carga, eficiencia operativa y riesgo/penalizaciones (Monterrey–Laredo). Elaboración propia.}
\end{figure}
El factorial mixto 9×3 confirma interacción significativa par O-D × peso de rentabilidad (p < 0,0001): la respuesta depende del corredor.

\section{Discusión y conclusiones}
La variabilidad de formaciones líderes por par O-D confirma la necesidad de enfoques MCDM en logística ferroviaria [10, 18]: la capa de carga incorpora sinergias de tonelaje, diversificación de riesgo y penalizaciones HM ausentes en políticas mono-vagón. La estabilidad 95 \% (sensibilidad) y concordancia 4/5 (validación cruzada) en Monterrey–Laredo, junto con la capa de ruteo sobre alternativas Dijkstra, refuerzan la utilidad operativa del marco bicapa.

El ANOVA factorial 33 aporta evidencia de efectos no aditivos entre rentabilidad/carga y eficiencia operativa, con curvatura significativa. El factorial mixto confirma que la recomendación es contingente al corredor. El trabajo extiende evaluaciones TOPSIS ferroviarias [12, 15] al nivel de composición multi-vagón, complementando estudios macro mexicanos [4, 7]. Entre las limitaciones destacan: (i) datos sintéticos que ilustran el método sin predicir desempeño operativo real; (ii) espacio combinatorio acotado (top-12, máx. 80 formaciones) y pesos fijos de capa de carga; (iii) inferencia ANOVA interna al escenario simulado. El simulador del norte de México funciona como laboratorio de decisión para operadores y docencia; líneas futuras incluyen calibración con datos Ferromex/CPKC, ponderación por entropía o AHP, y extensión del diseño factorial a otros pares O-D.

En síntesis:

  \item Se diseñó un marco TOPSIS bicapa (ruteo + carga) con pre-filtro interno para formaciones de 1–6 vagones en el corredor norte de México.
  \item Nueve pares O-D exhiben líderes predominantemente de tres vagones ($C_i$ 0.647–0.785); la composición óptima depende del par y del subconjunto factible.
  \item Monterrey–Laredo muestra estabilidad 95 \% ante sensibilidad ±20 \% y concordancia 4/5 en validación cruzada.
  \item El diseño factorial 33 confirma efectos significativos de rentabilidad/carga, eficiencia operativa, interacción AB y curvatura; el factorial mixto valida contingencia por corredor.
  \item Los materiales de apoyo publicados en https://github.com/Selenofilia/OptimizaciOn-Multicriterio-de-Trayecto-y-Estrategias-de-Carga-en-Transporte-Ferroviario- (texto del artículo, diseño experimental y resultados numéricos) permiten interpretar rankings, validación y ANOVA del estudio.
\subsubsection*{}
\textbf{Agradecimientos. }Los autores agradecen a la Universidad Autónoma de Guerrero y a la Universidad Autónoma de Ciudad Juárez el apoyo institucional brindado para el desarrollo de esta investigación. Asimismo, reconocen a Joa Pombo (ORCID: https://orcid.org/0000-0002-5877-1989) por la provisión del dataset comercial utilizado en este estudio.

\subsubsection*{}
\textbf{Declaración de intereses. }Los autores declaran no tener conflictos de interés relevantes para el contenido de este artículo.

\subsubsection*{}
\textbf{Declaración sobre inteligencia artificial. }Los autores utilizaron asistentes de IA generativa para borradores y rutinas de análisis; revisaron y asumen responsabilidad del contenido final.

\paragraph{}
\textbf{Disponibilidad de software y datos. }El texto del artículo (paper.tex), el diseño experimental, la red ferroviaria (data/cities.json) y los resultados numéricos (data/paper\_results.json, data/factorial\_experiment.json) están disponibles en https://github.com/Selenofilia/OptimizaciOn-Multicriterio-de-Trayecto-y-Estrategias-de-Carga-en-Transporte-Ferroviario-, versión v1.0-paper. El experimento reportado emplea el subconjunto núcleo de 587 estrategias.

\section{Referencias}
  \bibitem{Behzadian}
  Behzadian, M., Otaghsara, S.K., Yazdani, M., Ignatius, J.: A state-of-the-art survey of TOPSIS applications. Expert Syst. Appl. 39(17), 13051–13069 (2012). https://doi.org/10.1016/j.eswa.2012.05.056
  \bibitem{Blagojevi}
  Blagojević, A., Stević, Ž., Marinković, D., Kasalica, S., Rajilić, S.: A novel entropy-fuzzy PIPRECIA-DEA model for safety evaluation of railway traffic. Symmetry 12(9), Article 1479 (2020). https://doi.org/10.3390/sym12091479
  \bibitem{Broniewicz}
  Broniewicz, E., Ogrodnik, K.: Application potential of MCDM/MCDA methods in transport — Literature review and case study. Sustainability 17(17), Article 7671 (2025). https://doi.org/10.3390/su17177671
  \bibitem{CedilloCampos}
  Cedillo-Campos, M.G., García-Ortega, M.G., Martner-Peyrelongue, C.D., Saucedo-Martínez, J.A., Ponce-Ceja, N.: Flujos de carga automotriz y su impacto en la infraestructura ferroviaria en México: Un enfoque de fluidez en la cadena de suministro. Ing. Investig. Tecnol. 18(1), 87–99 (2017). https://doi.org/10.22201/fi.25940732e.2017.18n1.008
  \bibitem{Chakraborty}
  Chakraborty, S.: TOPSIS and modified TOPSIS: A comparative analysis. Decis. Anal. J. 2, Article 100021 (2022). https://doi.org/10.1016/j.dajour.2021.100021
  \bibitem{Dijkstra}
  Dijkstra, E.W.: A note on two problems in connexion with graphs. Numer. Math. 1(1), 269–271 (1959). https://doi.org/10.1007/BF01386390
  \bibitem{GarcaOrtega}
  García-Ortega, M.G., Morales Pérez, M.C.G.: Características de los corredores y flujos ferroviarios de carga intermodal en México, datos 2017 (Publicación Técnica No. 611). Instituto Mexicano del Transporte (2020). https://trid.trb.org/View/1763261
  \bibitem{Huang}
  Huang, W., Shuai, B., Sun, Y., Wang, Y., Antwi, E.: Using entropy-TOPSIS method to evaluate urban rail transit system operation performance: The China case. Transp. Res. Part A 111, 292–303 (2018). https://doi.org/10.1016/j.tra.2018.03.025
  \bibitem{Hwang}
  Hwang, C.-L., Yoon, K.: Multiple Attribute Decision Making: Methods and Applications, a State-of-the-Art Survey. Springer, Heidelberg (1981). https://doi.org/10.1007/978-3-642-48318-9
  \bibitem{Jung}
  Jung, H., Kim, J., Shin, K.: Importance analysis of decision-making factors for selecting international freight transportation mode. Asian J. Shipp. Logist. 35(1), 55–62 (2019). https://doi.org/10.1016/j.ajsl.2019.03.008
  \bibitem{Mantovani}
  Mantovani, S., Morganti, G., Umang, N., Crainic, T.G., Frejinger, E., Larsen, E.: The load planning problem for double-stack intermodal trains. Eur. J. Oper. Res. 267(1), 107–119 (2018). https://doi.org/10.1016/j.ejor.2017.10.051
  \bibitem{Marchetti}
  Marchetti, D., Wanke, P.: Efficiency of the rail sections in Brazilian railway system, using TOPSIS and a genetic algorithm to analyse optimized scenarios. Transp. Res. Part E 135, Article 101858 (2020). https://doi.org/10.1016/j.tre.2020.101858
  \bibitem{Meethom}
  Meethom, W., Koohathongsumrit, N.: An integrated potential assessment criteria and TOPSIS based decision support system for road freight transportation routing. Appl. Sci. Eng. Prog. 13(4), 312–326 (2020). https://doi.org/10.14416/j.ijast.2019.01.002
  \bibitem{Pandey}
  Pandey, V., Komal, Dincer, H.: A review on TOPSIS method and its extensions for different applications with recent development. Soft Comput. 27(23), 18011–18039 (2023). https://doi.org/10.1007/s00500-023-09011-0
  \bibitem{Petrovi}
  Petrović, N., Mihajlović, J., Jovanović, V., Ćirić, D., Živojinović, T.: Evaluating annual operation performance of Serbian railway system by using multiple criteria decision-making technique. Acta Polytech. Hung. 20(1), 157–168 (2023). https://doi.org/10.12700/aph.20.1.2023.20.11
  \bibitem{Rahman}
  Rahman, M.: Choosing multimodal freight mix: An integrated multi-objective multicriteria approach. J. Transp. Technol. 14(3), 402–422 (2024). https://doi.org/10.4236/jtts.2024.143024
  \bibitem{SteadieSeifi}
  SteadieSeifi, M., Dellaert, N.P., Nuijten, W., Van Woensel, T., Raoufi, R.: Multimodal freight transportation planning: A literature review. Eur. J. Oper. Res. 233(1), 1–15 (2014). https://doi.org/10.1016/j.ejor.2013.06.055
  \bibitem{Thompson}
  Thompson, E., Abudu, R., Zheng, S.: Empirical analysis of multiple-criteria decision-making (MCDM) process for freight transportation mode selection. J. Transp. Technol. 12(1), 28–41 (2022). https://doi.org/10.4236/jtts.2022.121002
\end{thebibliography}

\end{document}
